\section{Benchmark Construction}

\paragraph{Released data.}
We use the released SeekUI/VSGUI subset containing 1,362 present examples, 646 unique images, and 850 unique targets. The data audit shows that this released subset is text-target dominated: all target identifiers use the \texttt{txt} prefix, with only two examples having empty target text. This makes the current benchmark appropriate for studying text-target absence, while non-text and icon-like targets remain a follow-up direction.

\paragraph{Synthetic absent stress test.}
To evaluate target absence at scale, we construct a balanced synthetic present/absent benchmark. The 1,362 original examples are retained as present examples. We then construct 1,362 absent examples by pairing screenshots with target cues that should not appear on the screen, yielding 2,724 total examples. This stress test is designed to expose forced-choice behavior: the screenshot is still realistic, and the query is still plausible, but the model should not be forced to ground it.

\paragraph{Sanity checks.}
Because synthetic absence can introduce artifacts, we audit the benchmark before treating results as meaningful. The audit checks missing images, duplicate examples, annotation conflicts, target text mismatches, and OCR leakage. Annotation conflicts are rare, while OCR leakage is nontrivial; therefore OCR-heavy results must be interpreted with care. We also found that random train/test splitting leaks images, so the image-disjoint split is the cleaner held-out evaluation and the primary result in this draft.

\paragraph{Realistic absent validation.}
To reduce reliance on synthetic absence, we additionally create a small manually reviewed validation set with 100 examples: 50 present rows and 50 realistic absent queries. The absent queries are plausible user goals for the corresponding screen and manually checked to be absent. This set is not large enough to replace the synthetic benchmark, but it provides an initial external-validity check.
